% Generated draft: requires mathematical and source review.
\documentclass[11pt]{article}
\usepackage[margin=25mm]{geometry}
\usepackage{fontspec}
\setmainfont{Latin Modern Roman}
\usepackage{amsmath,amssymb,mathtools,unicode-math}
\setmathfont{Latin Modern Math}
\usepackage{enumitem,needspace,xurl,hyperref}
\setlength{\parindent}{0pt}
\setlength{\parskip}{5pt}
\setlength{\emergencystretch}{3em}
\newcommand{\N}{\mathbb{N}}
\newcommand{\Z}{\mathbb{Z}}
\newcommand{\Q}{\mathbb{Q}}
\newcommand{\R}{\mathbb{R}}
\newcommand{\C}{\mathbb{C}}
\newcommand{\F}{\mathbb{F}}
\newcommand{\A}{\mathbb{A}}
\newcommand{\PP}{\mathbb{P}}
\newcommand{\E}{\mathbb{E}}
\newcommand{\eps}{\varepsilon}
\newcommand{\dd}{\,\mathrm{d}}
\begin{document}
\section*{UnsolvedMath ID: 8800088}
\textbf{Source problem code:} catalogue entry
\subsection{Definitions and mathematical statement}
For a smooth projective geometrically connected curve $C/\F_q$, the degree-zero divisor class group is
\[
\operatorname{Cl}^0(C)=\{\text{$\F_q$-rational divisors of degree zero}\}/\{\text{principal divisors}\}.
\]
It is a finite abelian group. Its discrete logarithm problem, for a specified cyclic subgroup, gives classes $A,B$ with $B\in\langle A\rangle$ and asks for $x$ modulo $\operatorname{ord}(A)$ such that $B=xA$. The full class group need not be cyclic.

The source asks for efficient reductions in both directions between that problem and an infrastructure discrete logarithm problem. Arithmetic infrastructures typically use reduced ideal/divisor representatives together with distance information, and operations adapted to those representatives; they are not automatically groups with the same law as a class group. The exact infrastructure, distance space and logarithm task are not defined in the record.

A two-way reduction target must give effective transformations and oracle calls converting each specified task into the other while controlling input size, auxiliary information, precision and running time. A common asymptotic cost claim for two separate algorithms is not such a reduction. The source does not choose many-one versus oracle reductions, deterministic versus randomized guarantees, preprocessing, or whether quantities such as the regulator are known.

The summary mentions number-field and function-field settings, including quadratic and higher-degree cases, whereas the displayed divisor definition is the standard finite-field-curve instance. Those settings have different degree, distance and finiteness conventions and cannot be merged without explicit data. Thus the requested two-direction comparison is retained as a programme awaiting its field/infrastructure scope.
\subsection{Short English statement}
Convert divisor-class and infrastructure discrete logarithm instances efficiently in each direction. Specify the arithmetic setting, infrastructure distances, available regulator data and reduction model.
\subsection{Sources}
\begin{itemize}
\item Local UnsolvedMath registry, permanent ID 8800088, exact source code catalogue entry, ``Converting between divisor-class and infrastructure discrete logarithms''; attributed to Renate Scheidler, no proposed year supplied. Registry status partially solved, classification PARTIAL-PROGRESS.
\item \emph{IPAM Workshop on Number Theory and Cryptography---Open Problems} (2006), attachment 6665, PDF pages 184--188, printed slide 29, \url{https://www.ipam.ucla.edu/programs/workshops/workshop-i-number-theory-and-cryptography-open-problems/?tab=schedule}. Supplied pdf-and-legacy-ppt visual audit accessed July 29, 2026; provenance records 17 PDFs/two PPT decks inspected with repeated builds excluded. No PDF/PPT was accessed here.
\item The supplied summary names Scheidler, Jacobson, Teske and Galbraith--Menezes for special-setting comparisons in number fields and function fields, without exact references, hypotheses or reductions. It reports no verified general efficient two-way equivalence.
\item Provenance flags current primary-literature status and redistribution terms NEEDS\_REVIEW.
\end{itemize}
\end{document}
